% 灵敏度分析
\chapter{灵敏度分析：$B^*$ 对期待与约束的响应}\label{chap:sensitivity}

第~\ref{chap:overview}章把 $B^*$ 刻画为期待与约束的函数：$B^* = f(\text{期待},\ \text{约束})$，
并在给定技术选型下用``固定 $B$ 的可行性 LP + 二分搜索''求得最大端口无阻塞带宽。
本章回答紧随其后的问题：如果期待或约束发生微小变化，$B^*$ 会如何响应？
这个响应即\textbf{灵敏度}，它把``可行/不可行''的判定提升为``往哪个方向改、改多少、收益多少''的设计指导。

\section{目标：$\partial B^*/\partial\theta$}\label{sec:sens-objective}

灵敏度分析要回答的不是``哪个约束绑定''，而是``$B^*$ 对参数 $\theta$ 的梯度''。
最直接的实例是约束右手边：$\partial B^*/\partial b_i$ = 约束 $i$ 的 rhs 松弛 1 个单位、$B^*$ 增加多少。
这是架构师真正想要的量：投资哪个约束（+1 $\mu$bump / +1\,W / +1 lane）换多少带宽。

更一般的 $\theta$ 覆盖两类旋钮（见第~\ref{sec:sens-knobs}节）：
\textbf{约束旋钮}（物理参数 $p,\eta,R_{\text{th}},\alpha_d,\beta_P$ 等）与
\textbf{期待旋钮}（无阻塞要求的严格程度——哪些流量模式进 $\mathcal{R}$ 集合）。
本章把二者统一为对参数 $\theta$ 的梯度 $\partial B^*/\partial\theta$。

\section{数学基础}\label{sec:sens-math}

\subsection{参数化优化与包络定理}\label{sec:sens-envelope}

把 $B^*$ 视为参数化优化的最优值。为表述简洁，先考虑 rhs 独立于 $B$ 的基准情形，
即约束可写成 $B\cdot(A_i\cdot\mathbf{L})\le b_i$（$A_i$ 为约束 $i$ 消元后对 $\mathbf{L}$ 的系数行，$b_i$ 为常数 rhs）：

\begin{equation}\label{eq:sens-param}
    B^*(b) \;=\; \max_{B,\ \mathbf{L}\in\mathcal{F}}\; B
    \qquad\text{s.t.}\qquad
    B \cdot (A_i \cdot \mathbf{L}) \;\le\; b_i,\quad \forall i
\end{equation}

其中 $\mathcal{F}$ 是性能侧（路由 + 定义式，式~\eqref{eq:route-flow}）决定的 $\mathbf{L}$ 可行域，与 $b$ 无关。
拉格朗日 $\mathcal{L}=B+\sum_i\lambda_i\,(b_i-B\cdot(A_i\cdot\mathbf{L}))$。
\textbf{包络定理}给出 $\partial B^*/\partial b_i=\lambda_i^*$——约束 $i$ 的 KKT 乘子即带宽影子价格。
对 $B$ 求偏导得 KKT 条件：

\begin{equation}\label{eq:sens-kkt}
    1 - \sum_i \lambda_i\,(A_i \cdot \mathbf{L}^*) = 0
    \qquad\Longrightarrow\qquad
    \sum_i \lambda_i\,(A_i \cdot \mathbf{L}^*) = 1
\end{equation}

互补松弛：$\lambda_i\ge 0$，$\lambda_i\,(b_i-B^*\cdot(A_i\cdot\mathbf{L}^*))=0$。

\subsection{非凸问题的处理：二分 + 事后闭式}\label{sec:sens-posthoc}

注意 $B\cdot(A_i\cdot\mathbf{L})$ 是双线性项，$\max_{B,\mathbf{L}}B$ 受此类约束是\textbf{非凸}问题——
KKT 只是驻点必要条件，不保证全局最优。因此本文的 $B^*$ \textbf{不}通过求解式~\eqref{eq:sens-param}得到，
而是按第~\ref{sec:ov-unified}节的``固定 $B$ 可行性 LP + 二分搜索''求得：
每次固定 $B$，约束是 $\mathbf{L}$ 的线性不等式（凸 LP）；
$B^*$ 的全局性由``可行性关于 $B$ 单调 + 二分''保证，不依赖 KKT。

$\lambda$ 的取得是\textbf{事后闭式}：二分收敛后取最后一次可行解 $\mathbf{L}^*$，
对每个绑定约束 $j$ 用

\begin{equation}\label{eq:sens-lambda}
    \lambda_j = \frac{1}{A_j \cdot \mathbf{L}^*}
\end{equation}

（因绑定处 $B^*=b_j/(A_j\cdot\mathbf{L}^*)$ 刻画同一超平面）。
该闭式值与 $\max B$ 问题在 $B^*$ 处的 KKT 乘子数值一致，但不依赖求解非凸 $\max B$ 问题。

式~\eqref{eq:sens-lambda} 与第~\ref{sec:sens-knobs}节的 $\partial B^*/\partial\theta=\sum_i\lambda_i\,\partial g_i/\partial\theta$
挂 \textbf{Milgrom--Segal (2002) 包络定理}：它只要求 value function 在参数处可微 + 约束规范，
\textbf{不要求凸性}；在 $B^*$ 可微点成立，两个约束同时绑定的退化点只给方向导数/次梯度（见下）。

\subsection{单约束绑定与退化情况}\label{sec:sens-degenerate}

若仅约束 $j$ 绑定（其余 slack，$\lambda_i=0$），式~\eqref{eq:sens-kkt} 直接给出

\begin{equation}\label{eq:sens-single}
    \lambda_j = \frac{1}{A_j \cdot \mathbf{L}^*},
    \qquad
    \frac{\partial B^*}{\partial b_j} = \frac{1}{A_j \cdot \mathbf{L}^*}
\end{equation}

物理含义：$A_j\cdot\mathbf{L}^*$ 是``约束 $j$ 单位带宽的 lhs 占用''，其倒数是``每松弛 1 单位 rhs 换多少 Gbps''。

若约束集 $\mathcal{J}$ 同时绑定，式~\eqref{eq:sens-kkt} 只给出
$\{\lambda_j\ge 0,\ \sum_{j\in\mathcal{J}}\lambda_j\,(A_j\cdot\mathbf{L}^*)=1\}$：
单独松弛一个约束 $B^*$ 不增（其余仍挡路），$\lambda_j$ 需联合确定，且退化不唯一。
此时``先改哪个''由绑定演化序列（沿 $B$ 轴谁先绑定）回答。

（实现注记：二分搜索每次求解的``固定 $B$ 可行性判定''无目标函数，
其对偶不可靠，且含义是 $\partial(\min\sum_e L_e)/\partial b_i$（固定 $B$ 下），
与 $\partial B^*/\partial b_i$ 不是同一个量；故 $\lambda$ 一律用式~\eqref{eq:sens-lambda} 事后闭式求取，不依赖可行性对偶。）

\section{两个旋钮：约束旋钮与期待旋钮}\label{sec:sens-knobs}

接入 die 缩放模型（式~\eqref{eq:die-scale}）后，约束 rhs 本身是 $B$ 的函数
（$\mathbf{N}^{\text{total}}(B)$、$\mathbf{P}_{peak}(B)$），``$b_i$ 独立于 $B$''的前提失效。
更本质的灵敏度目标是对\textbf{物理参数} $\theta$ 求导：

\begin{equation}\label{eq:sens-dtheta}
    \frac{\partial B^*}{\partial \theta}
    \;=\; \sum_i \lambda_i^* \cdot \frac{\partial g_i}{\partial \theta}
\end{equation}

其中约束 $i$ 记为 $g_i(B,\mathbf{L},\theta)\le 0$，$\lambda_i^*$ 是 KKT 乘子（单约束闭式式~\eqref{eq:sens-single}），
$\partial g_i/\partial\theta$ 是约束对参数的显式偏导（对第~\ref{chap:overview}章线性化系数的显式求导）。

统一两种情况：

\begin{itemize}
\item \textbf{约束旋钮}：$\theta$ = 物理参数——bump pitch $p$、阵列利用率 $\eta$、热阻 $R_{\text{th}}$、
die 边长 $d_0$、缩放斜率 $\alpha_d$、功耗斜率 $\beta_P$，用式~\eqref{eq:sens-dtheta} 的通用形式。
$\theta$ = rhs 常数 $b_i$ 是特例（$\partial g_i/\partial b_i=1$，回到 $\partial B^*/\partial b_i=\lambda_i$）。
\item \textbf{期待旋钮}：$\theta$ = $\mathcal{R}$ 集合的严格程度——$\mathcal{R}$ 变小时需求下界 $\mathbf{L}$ 变小，
经同一梯度结构传至 $B^*$：期待降低（少要求一些流量模式无阻塞），$B^*$ 提高。
\end{itemize}

两个旋钮方向明确：\textbf{约束放宽}（工艺、散热、功耗调度改善）或\textbf{期待降低}，$B^*$ 都提高；反之亦然。

\section{归一化与可比性}\label{sec:sens-norm}

不同约束 rhs 单位不同（bump 数 / W / lane 数），直接比较影子价格无意义。
归一化后给架构师可比的投资优先级：每 +1 $\mu$bump $\to X$\,Gbps；每 +1\,W 散热 $\to Y$\,Gbps；
比值 $\nu_{\text{bump}}/\nu_{\text{thermal}}$ 判断投 bump pitch 缩减还是更激进的冷却。
由此，灵敏度把``可行/不可行''的判定提升为``往哪个方向改、改多少、收益多少''的设计指导。
